\documentclass[10pt,a4paper]{amsart}
\usepackage[margin=19mm]{geometry}
\usepackage{amsmath,amssymb,mathtools}
\usepackage[scr=rsfs,cal=euler]{mathalfa}
\usepackage{xfrac}
\usepackage[unicode,hidelinks,bookmarks=false]{hyperref}
\newcommand{\ie}{i.e.\ }
\newcommand{\cf}{cf.\ }
\newcommand{\resp}{respectively\ }
\newcommand{\Subsection}{\subsection}
\providecommand{\nobreakdash}{\nobreak\hbox{-}}
\setlength{\emergencystretch}{2em}
\multlinegap0pt
\usepackage{amssymb}
\newtheorem{lemma}{Lemma}
\title{Three Graffiti.pc Conjectures on Largest Induced Trees:\\ Proofs of Conjectures 141, 142, and 143}
\author{Alper Ferudun}
\date{Source: arXiv:2608.01396v1 (CC BY 4.0)}
\begin{document}
\maketitle

\noindent\emph{Source and licence.} Three Graffiti.pc Conjectures on Largest Induced Trees:\\ Proofs of Conjectures 141, 142, and 143. Version arXiv:2608.01396v1 (CC BY 4.0), distributed by arXiv under the Creative Commons Attribution 4.0 International licence, http://creativecommons.org/licenses/by/4.0/, as recorded in the arXiv licence metadata for this exact version and shown on the abstract page https://arxiv.org/abs/2608.01396v1. Full licence notice: \texttt{licenses/arxiv-2608.01396-CC-BY-4.0.txt}.\par\smallskip

\noindent\textit{Editorial scope.} Counted unit: the article's lemma w142-lem:rooted-cycle (source0.tex lines 453--464). The article's complete original proof of it is delivered with it (source0.tex lines 466--578, one \texttt{proof} environment of 113 lines). No mathematics is written, repaired or re-derived: the statement and the proof are the article's own lines, in the article's own notation, and no retained line is substituted or re-flowed.  Retained uncounted as the source-local context the proof needs: lines 436--440.  Nothing was omitted from any retained span, so the package carries no omission record.  No retained line contains a citation, so the wrapper carries no bibliography and no bibliography entry.  No retained line contains a figure environment, an image-inclusion command or drawing code, so no image is delivered and the package declares no asset dependency.  Editorial formatting only, in addition: an AMS \texttt{amsart} preamble that restates the article's own declarations and macros used by the retained lines, namely the environments \texttt{lemma} on separate counters, together with the standard packages those lines need.  The retained environments print in this document as Lemma 1 in order of appearance, whereas the article prints the same items with the numbering of its own full document.  The complete native source of the article as distributed in the arXiv e-print for this exact version is bundled unchanged as \texttt{sources/206622/source0.tex}.
\begin{lemma}[Cycles in a minimal connector]\label{w142-lem:minimal-connector}
Let $J$ be connected and $Q\subseteq V(J)$.  If no proper induced connected
subgraph of $J$ contains $Q$, then every cycle $L$ of $J$ has
$|V(L)|\le |Q|$.
\end{lemma}

\begin{lemma}[Rooted shortest cycle]\label{w142-lem:rooted-cycle}
Assume $g\ge5$, and let $S\subseteq V(G)$ with $|S|\le3$.  There exist a
shortest cycle $K$, a vertex $z\in V(K)\setminus S$, and a set
$F\subseteq V(G)\setminus V(K)$ such that
\[
 S\subseteq (V(K)\setminus\{z\})\cup F,
\]
$G[F]$ is a forest, and every component $C$ of $G[F]$ satisfies
\[
 e(C,V(K)\setminus\{z\})=1.
\]
\end{lemma}

\begin{proof}
The empty case is immediate.  Fix an arbitrary shortest cycle $K$, and
choose $X\supseteq V(K)\cup S$ of minimum cardinality such that
$H:=G[X]$ is connected.

Every component $C$ of $H-V(K)$ contains a terminal.  Indeed, it has an
edge to $K$, since $H$ is connected; if $S_C:=S\cap V(C)$ were empty, the
whole component could be deleted while preserving connectedness and all
required vertices.  Put
\[
 A_C:=\{a\in V(C):N_G(a)\cap V(K)\ne\varnothing\}.
\]
Form $J_C$ from $G[C]$ by adjoining a new vertex $r_C$ adjacent precisely
to $A_C$.  This graph is vertex-minimal among induced connected subgraphs
containing $S_C\cup\{r_C\}$.  Otherwise, replacing $C$ by a smaller
connector would leave every resulting component attached to $K$ and would
contradict the choice of $X$.

Lemma~\ref{w142-lem:minimal-connector} shows that every cycle in $J_C$ has length
at most $|S_C|+1\le4$.  A cycle contained in $G[C]$ would also be a cycle of
$G$ and would have length at least $g\ge5$.  Consequently
\begin{equation}\label{w142-eq:C-tree}
  G[C]\text{ is a tree}.
\end{equation}

Let $R_C$ be the minimal subtree of $C$ containing $A_C$.  We claim
\begin{equation}\label{w142-eq:R-small}
  |V(R_C)|\le |S_C|.
\end{equation}
For $v\in V(R_C)$, let $Q_v$ be the component of $J_C-v$ containing $r_C$.
Minimality supplies $s_v\in S_C\setminus V(Q_v)$.  For $s\in S_C$, let
$c(s)$ be the first vertex of $R_C$ on the unique path from $s$ to $R_C$.
Every nonempty component of $R_C-v$ contains a vertex of $A_C$, by the
minimality of $R_C$.  Thus, if $v\ne c(s)$, the path from $s$ to $c(s)$ and
then to such an attachment gives an $s$--$r_C$ path in $J_C-v$; hence
$s\in Q_v$.  Applying this to $s_v$ gives $c(s_v)=v$.  The map
$v\mapsto s_v$ is injective, proving~\eqref{w142-eq:R-small}.

Every $a\in A_C$ has exactly one neighbour on $K$.  Existence follows from
the definition, while two distinct neighbours and their shorter $K$-arc
would form a cycle of length at most $\lfloor g/2\rfloor+2<g$.  Denote the
unique root by $\rho(a)$.  If $a,a'\in A_C$ are distinct, then
\begin{equation}\label{w142-eq:attachment-distance}
 d_C(a,a')\le |V(R_C)|-1\le2
\end{equation}
and their roots are distinct: equal roots, together with the
$a$--$a'$ path in $C$, would form a cycle of length at most four.  Moreover,
the two root edges, that path, and a shorter arc of $K$ give
\begin{equation}\label{w142-eq:girth-roots}
 d_K(\rho(a),\rho(a'))+d_C(a,a')+2\ge g.
\end{equation}

These inequalities sharply restrict the number of attachments.  If
$g\ge9$, then two attachments would make the left side of
\eqref{w142-eq:girth-roots} at most $\lfloor g/2\rfloor+4<g$, so there is only
one.  For $g\in\{7,8\}$, two attachments cannot be adjacent in $C$; three
attachments would fill the three vertices of $R_C$ and contain an adjacent
pair.  For $g=6$, if $R_C=p-q-r$ consisted of three attachments, applying
\eqref{w142-eq:girth-roots} to $p,q$ and to $q,r$ would force both $\rho(p)$ and
$\rho(r)$ to be the unique antipode of $\rho(q)$ on the $6$-cycle,
contrary to distinctness.  Thus every component has at most two
attachments, except that for $g=5$ one component may have three.  Also, at
most one component has two or more attachments, because such a component
contains at least two terminals and $|S|\le3$.

Suppose first that every component has one attachment $a_C$.  The set
\[
 Q=(S\cap V(K))\cup\{\rho(a_C):C\text{ a component of }H-V(K)\}
\]
has cardinality at most $|S|\le3$.  Since $g\ge5$, choose $z\in V(K)\setminus
Q$ and take $F=X\setminus V(K)$.  Equation~\eqref{w142-eq:C-tree} gives a forest,
each component has its single attachment in $Q$, and $z\notin S$.

Next suppose that one component $C_0$ has exactly two attachments with roots
$u,v$.  It contains at least two terminals, so the set
\[
 Q=(S\cap V(K))\cup
   \{\rho(a_C):C\ne C_0\}
\]
has size at most one.  Choose $z\in\{u,v\}\setminus Q$ and again put
$F=X\setminus V(K)$.  Every ordinary component has one edge into $K-z$,
while $C_0$ has exactly the two root edges and loses exactly one of them.
Again $z\notin S$.

It remains to handle three attachments.  Necessarily $g=5$,
$|V(R_C)|=|S_C|=|S|=3$, there are no other off-cycle components, and
$R_C=p-q-r$ with all three vertices attachments.  Write
\[
 u=\rho(p),\qquad v=\rho(q),\qquad w=\rho(r).
\]
Equation~\eqref{w142-eq:girth-roots} gives $d_K(u,v)=d_K(v,w)=2$.  In cyclic
order, write $K=u\alpha v\beta w u$.  Then
\[
  K'=u\alpha v q p u
\]
is another shortest $5$-cycle.

The injection used to prove~\eqref{w142-eq:R-small} is now a bijection.  Hence
there are distinct terminals $s_p,s_q,s_r$ with $c(s_t)=t$.  Let $L_t$ be
the unique $t$--$s_t$ path in the tree $C$, and set
\[
 F'=(V(L_p)\setminus\{p\})\cup
    (V(L_q)\setminus\{q\})\cup V(L_r),
 \qquad z=\alpha.
\]
The three displayed path-pieces are pairwise disjoint and anticomplete, for
an edge between two of them would create a cycle in the tree $C$.  The first
two nonempty pieces attach once to $K'$ at $p$ and $q$, respectively.  The
third attaches once through the edge $rq$; its other old-cycle edge $rw$
does not meet $K'$.  Thus $F'$ is admissible for $(K',z)$.  Finally each
$s_t$ belongs either to its path-piece or, when $s_p=p$ or $s_q=q$, directly
to $K'-z$.  This completes the exceptional case and the proof.
\end{proof}

\end{document}
